\documentclass[11pt]{article}

\usepackage[T1]{fontenc}
\usepackage[utf8]{inputenc}
\usepackage{lmodern}
\usepackage[margin=0.75in]{geometry}
\usepackage{xcolor}
\usepackage{booktabs}
\usepackage{tabularx}
\usepackage{array}
\usepackage{listings}
\usepackage{hyperref}
\usepackage{titlesec}
\usepackage{enumitem}
\usepackage[most]{tcolorbox}
\usepackage{microtype}

\definecolor{accent}{HTML}{0B6E69}
\definecolor{accentdark}{HTML}{084C49}
\definecolor{muted}{HTML}{5B6770}
\definecolor{soft}{HTML}{F2F7F6}
\definecolor{line}{HTML}{C9D8D5}
\definecolor{codebg}{HTML}{F7F9FA}

\hypersetup{
  colorlinks=true,
  linkcolor=accentdark,
  urlcolor=accentdark,
  citecolor=accentdark,
  pdftitle={RTX 5090 Low-Precision Product-Pattern Width Report},
  pdfauthor={tensor-cores-numerical-behavior}
}

\titleformat{\section}
  {\Large\bfseries\color{accentdark}}
  {\thesection}{0.55em}{}
\titleformat{\subsection}
  {\large\bfseries\color{accentdark}}
  {\thesubsection}{0.55em}{}
\titlespacing*{\section}{0pt}{1.2em}{0.5em}
\titlespacing*{\subsection}{0pt}{0.85em}{0.3em}

\setlist[itemize]{leftmargin=1.2em, itemsep=0.2em, topsep=0.25em}
\setlength{\parindent}{0pt}
\setlength{\parskip}{0.5em}
\renewcommand{\arraystretch}{1.18}

\lstdefinestyle{reportcode}{
  basicstyle=\ttfamily\footnotesize,
  backgroundcolor=\color{codebg},
  frame=single,
  rulecolor=\color{line},
  framerule=0.5pt,
  breaklines=true,
  columns=fullflexible,
  keepspaces=true,
  showstringspaces=false,
  tabsize=2
}
\lstset{style=reportcode}

\newcolumntype{Y}{>{\raggedright\arraybackslash}X}
\newcolumntype{R}{>{\raggedleft\arraybackslash}p{0.78in}}
\newcolumntype{C}{>{\centering\arraybackslash}p{0.82in}}

\newtcolorbox{summarybox}{
  enhanced,
  colback=soft,
  colframe=accent,
  boxrule=0.7pt,
  arc=2mm,
  left=8pt,
  right=8pt,
  top=8pt,
  bottom=8pt
}

\begin{document}
\pagenumbering{gobble}

\begin{titlepage}
\vspace*{0.3in}
{\Huge\bfseries\color{accentdark} RTX 5090 Low-Precision\\ Product-Pattern Width Report\par}
\vspace{0.18in}
{\Large Direct PTX \texttt{M - M + epsilon} product probes\par}
\vspace{0.22in}
{\color{muted}Date: 2026-05-09\par}
\vfill
\begin{summarybox}
\textbf{Headline finding.}
The earlier A-only pattern understated FP8 E4M3 range. With independent
\texttt{A*B} products, FP8 E4M3 x E4M3 reaches an extreme test width of
\textbf{36 bits}. That extreme case does not survive, and product-exponent
scanning finds \textbf{27 effective bits} for all tested FP8 variants. FP6 and
FP4 remain format-limited in this product-pattern probe.
\end{summarybox}
\vspace{0.22in}
\begin{tabularx}{\linewidth}{>{\bfseries}l Y}
\toprule
GPU & NVIDIA GeForce RTX 5090, compute capability 12.0 \\
Build target & \texttt{compute\_120a} / \texttt{sm\_120a} \\
Source & \path{src/tc_test_numerics-5090-lowp-reduction-pattern.cu} \\
Result files & \path{5090/result-5090-fp8-reduction-pattern.txt},
\path{5090/result-5090-fp6-reduction-pattern.txt},
\path{5090/result-5090-fp4-reduction-pattern.txt} \\
\bottomrule
\end{tabularx}
\end{titlepage}
\clearpage
\pagenumbering{arabic}

\section{Scope}

This report covers the standalone product-pattern probe listed on the title
page. It is separate from the accumulator-boundary and repeated-MMA probes. The
test uses direct inline PTX \texttt{mma.sync.aligned} instructions and does not
call cuBLASLt.

The probe reports externally visible behavior for selected product positions.
It does not claim to reconstruct the physical tensor-core reduction tree.

\section{Reading Guide}

This report uses three related but different quantities:

\begin{tabularx}{\linewidth}{l Y Y}
\toprule
Term & Meaning & Interpretation \\
\midrule
Extreme width & Gap between the largest finite A/B product and the smallest
positive A/B product. & Widest gap the format can generate in this test. \\
Max effective width & Widest product-exponent gap where epsilon remains visibly
retained after cancellation. & Measured retention boundary for this layout and
criterion. \\
Format-limited & Max effective width equals extreme width. & The test ran out
of representable product range before finding a failing boundary. \\
\bottomrule
\end{tabularx}

The key result is therefore: FP8 exposes a non-format-limited 27-bit retention
boundary in this selected layout; FP6 and FP4 remain capped by their
representable product ranges.

\section{Method}

The test maps independent product terms for output register \texttt{D0}. The
selected terms are:

\begin{tabularx}{\linewidth}{l Y R}
\toprule
Role & Position & Weight in D0 \\
\midrule
\texttt{+M} & \texttt{A0.byte0 * B0.byte0} & 4 \\
\texttt{-M} & \texttt{A0.byte1 * B0.byte1} & 4 \\
\texttt{epsilon} & \texttt{A0.byte2 * B0.byte2} & 4 \\
\bottomrule
\end{tabularx}

The cross terms among those positions are checked to be absent for
\texttt{D0}, so the pattern isolates three product contributions:

\[
  D0 = 4M - 4M + 4\epsilon
\]

The test enumerates every positive finite raw value in the A format and every
positive finite raw value in the B format. It forms all positive products,
groups them by product exponent, and scans exponent gaps:

\begin{summarybox}
\[
  effective\ width = exponent(M) - exponent(\epsilon) + 1
\]
\end{summarybox}

A case counts as retained when:

\begin{lstlisting}
observed > 0
abs(observed - expected) <= 0.5 * abs(expected)
\end{lstlisting}

The 50\% tolerance is intentional. This pattern detects whether the small term
is still visible after cancellation. It does not require exact arithmetic on
the retained term, so borderline FP8 cases are reported explicitly.

The sign pattern uses a negative A value for the \texttt{-M} term and keeps the
same B product magnitude. This keeps \texttt{+M} and \texttt{-M} symmetric
while allowing the third independent product to carry \(\epsilon\).

\section{Results}

\begin{tabularx}{\linewidth}{l Y R C R C}
\toprule
Family & Variant & Extreme width & Extreme survived & Max effective width & Format-limited \\
\midrule
FP8 & E4M3 x E4M3 & 36 & no & 27 & no \\
FP8 & E5M2 x E5M2 & 64 & no & 27 & no \\
FP8 & E4M3 x E5M2 & 50 & no & 27 & no \\
FP8 & E5M2 x E4M3 & 50 & no & 27 & no \\
FP6 & E2M3 x E2M3 & 12 & yes & 12 & yes \\
FP6 & E3M2 x E3M2 & 18 & yes & 18 & yes \\
FP4 & E2M1 x E2M1 & 8 & yes & 8 & yes \\
\bottomrule
\end{tabularx}

\subsection{Product Range}

\begin{tabularx}{\linewidth}{Y R R R R}
\toprule
Variant & Tested pairs & Survived & Max product & Min product \\
\midrule
FP8 E4M3 x E4M3 & 666 & 612 & 200704 & 3.8147e-06 \\
FP8 E5M2 x E5M2 & 2080 & 1332 & 3.2883e+09 & 2.3283e-10 \\
FP8 E4M3 x E5M2 & 1275 & 976 & 25690112 & 2.9802e-08 \\
FP8 E5M2 x E4M3 & 1275 & 976 & 25690112 & 2.9802e-08 \\
FP6 E2M3 x E2M3 & 78 & 78 & 56.25 & 0.015625 \\
FP6 E3M2 x E3M2 & 171 & 171 & 784 & 0.00390625 \\
FP4 E2M1 x E2M1 & 36 & 36 & 36 & 0.25 \\
\bottomrule
\end{tabularx}

\subsection{Survival Ratios}

\begin{tabularx}{\linewidth}{Y R Y}
\toprule
Variant & Survival ratio & Meaning \\
\midrule
FP8 E4M3 x E4M3 & 91.9\% & Wide gaps begin to fail; not format-limited. \\
FP8 E5M2 x E5M2 & 64.0\% & Much wider available range; many large gaps fail. \\
FP8 E4M3 x E5M2 & 76.5\% & Mixed range exposes the same 27-bit boundary. \\
FP8 E5M2 x E4M3 & 76.5\% & Same boundary in the opposite operand order. \\
FP6 E2M3 x E2M3 & 100\% & No failing case within representable product range. \\
FP6 E3M2 x E3M2 & 100\% & No failing case within representable product range. \\
FP4 E2M1 x E2M1 & 100\% & No failing case within representable product range. \\
\bottomrule
\end{tabularx}

\subsection{Extreme Cases}

\begin{tabularx}{\linewidth}{Y R R R C}
\toprule
Variant & Width & Observed & Expected & Survived \\
\midrule
FP8 E4M3 x E4M3 & 36 & 0 & 1.5259e-05 & no \\
FP8 E5M2 x E5M2 & 64 & 0 & 9.3132e-10 & no \\
FP8 E4M3 x E5M2 & 50 & 0 & 1.1921e-07 & no \\
FP8 E5M2 x E4M3 & 50 & 0 & 1.1921e-07 & no \\
FP6 E2M3 x E2M3 & 12 & 0.0625 & 0.0625 & yes \\
FP6 E3M2 x E3M2 & 18 & 0.015625 & 0.015625 & yes \\
FP4 E2M1 x E2M1 & 8 & 1 & 1 & yes \\
\bottomrule
\end{tabularx}

\subsection{Max-Width Retained Cases}

\begin{tabularx}{\linewidth}{Y R R R R R}
\toprule
Variant & Width & M exp & Eps exp & Observed & Expected \\
\midrule
FP8 E4M3 x E4M3 & 27 & 17 & -9 & 0.0078125 & 0.0153809 \\
FP8 E5M2 x E5M2 & 27 & 31 & 5 & 128 & 240 \\
FP8 E4M3 x E5M2 & 27 & 24 & -2 & 1 & 1.96875 \\
FP8 E5M2 x E4M3 & 27 & 24 & -2 & 1 & 1.96875 \\
FP6 E2M3 x E2M3 & 12 & 5 & -6 & 0.0625 & 0.0625 \\
FP6 E3M2 x E3M2 & 18 & 9 & -8 & 0.015625 & 0.015625 \\
FP4 E2M1 x E2M1 & 8 & 5 & -2 & 1 & 1 \\
\bottomrule
\end{tabularx}

The FP8 max-width retained cases are edge cases. The observed value is roughly
half of the expected epsilon contribution, but remains positive and within the
retention tolerance. Therefore, \textbf{27 bits} means the widest detectable
retained epsilon under this criterion, not exact 27-bit arithmetic.

The next wider gaps are not listed as pass cases because the test requires both
a positive result and a result close enough to the expected epsilon. This avoids
counting unrelated residue as a retained small term.

\section{Format Analysis}

\subsection{FP8 E4M3 x E4M3}

The earlier A-only pattern understated the available range. With A and B both
participating:

\begin{lstlisting}
max product = 200704
min product = 3.81469727e-06
product exponent gap = 17 - (-18) + 1 = 36
\end{lstlisting}

The extreme case is lost. The widest retained case is 27 bits, with
\(M\) at exponent 17 and \(\epsilon\) at exponent -9. This is not
format-limited.

The retained 27-bit case is near the tolerance edge:
\texttt{observed = 0.0078125} and
\texttt{expected = 0.0153808594}. This is a retention signal, not an exact
summation signal.

\subsection{FP8 E5M2 x E5M2}

E5M2 has a wider exponent range than E4M3, so the extreme product width reaches
64 bits. The extreme case is lost, and the widest retained case is again 27
bits. More representable product range does not increase the observed retained
width for this pattern.

The retained 27-bit case is also near the threshold:
\texttt{observed = 128} and \texttt{expected = 240}. E5M2 provides a larger
test range, but the widest retained product gap is not more accurate in
absolute terms.

\subsection{FP8 Mixed E4M3/E5M2}

Both mixed variants reach an extreme width of 50 bits and a max effective width
of 27 bits. The two directions match in this selected layout, but that should
not be generalized to every possible operand placement or output lane without
additional layout sweeps.

The mixed cases show that the 27-bit retained-width signal is not unique to a
same-format FP8 instruction under this selected independent-product layout.

\subsection{FP6 E2M3 x E2M3}

The entire tested product range is 12 bits, and the extreme case survives
exactly. This result is format-limited: it proves the tested path preserves the
full E2M3 product range used here, but it does not prove the internal retention
width is only 12 bits.

\subsection{FP6 E3M2 x E3M2}

The tested range expands to 18 bits because E3M2 has more exponent range than
E2M3. The extreme case survives exactly, so this is also format-limited. The
test does not find a failing FP6 E3M2 product-pattern case.

\subsection{FP4 E2M1 x E2M1}

The tested product range is 8 bits. The extreme case survives exactly. This is
strong evidence that the tested path can retain every product gap expressible
by FP4 E2M1 in this layout, but it cannot bound a wider internal tree because
the format cannot generate a wider product gap.

\section{Edge Cases and Limitations}

\begin{itemize}
\item The result is layout-sensitive. Other positions, lanes, or output
  registers may have different association behavior.
\item FP8's 27-bit retained cases pass because epsilon remains visible and
  within the 50\% retention tolerance, not because the result equals the
  expected value.
\item FP6 and FP4 are format-limited in this test. Their max effective width is
  the largest width generated by the product format here.
\item The grouping key is product exponent. Products in the same exponent
  bucket can have different significands.
\item This test covers direct \texttt{mma.sync.aligned.m16n8k32} PTX and one
  selected output register. It does not cover block-scaled MXFP instructions
  or full GEMM accuracy.
\end{itemize}

\section{Relationship To The Boundary Probe}

This report and \texttt{lowp-reduction-width-report.md} measure different
questions:

\begin{tabularx}{\linewidth}{l Y Y}
\toprule
Probe & Pattern & Reported signal \\
\midrule
Boundary scan & \texttt{D = C + 32} & When a full MMA update stops changing a
large external accumulator. \\
Product pattern & \texttt{M - M + epsilon} & Whether a tiny product term
survives internal cancellation. \\
\bottomrule
\end{tabularx}

The accumulator-boundary probe reports a 21-bit per-MMA update boundary. This
product-pattern probe reports FP8 retained cases up to 27 bits in a specific
cancellation layout. These values are complementary, not contradictory.

\section{Suggested Follow-Up Tests}

\begin{itemize}
\item Sweep additional independent product positions and output registers.
\item Tighten the retained-case threshold to separate epsilon visibility from
  approximate epsilon correctness.
\item Record the smallest failing gap above each max retained width.
\item Add block-scaled MXFP4/MXFP6 variants when direct PTX support is
  available.
\end{itemize}

\section{Verification}

\begin{lstlisting}[language=bash]
make -B test-5090-fp8-reduction-pattern \
        test-5090-fp6-reduction-pattern \
        test-5090-fp4-reduction-pattern

python3 run_tests.py -o 5090 \
        5090-fp8-reduction-pattern \
        5090-fp6-reduction-pattern \
        5090-fp4-reduction-pattern

compute-sanitizer --tool memcheck ./test-5090-fp8-reduction-pattern
compute-sanitizer --tool memcheck ./test-5090-fp6-reduction-pattern
compute-sanitizer --tool memcheck ./test-5090-fp4-reduction-pattern
\end{lstlisting}

The generated SASS contains direct \texttt{QMMA.16832} instructions:

\begin{tabularx}{\linewidth}{l Y}
\toprule
Target & Confirmed instruction family \\
\midrule
\texttt{test-5090-fp8-reduction-pattern} & \texttt{QMMA.16832.F32.E4M3/E5M2} \\
\texttt{test-5090-fp6-reduction-pattern} & \texttt{QMMA.16832.F32.E2M3/E3M2} \\
\texttt{test-5090-fp4-reduction-pattern} & \texttt{QMMA.16832.F32.E2M1} \\
\bottomrule
\end{tabularx}

\end{document}
